\documentclass[11pt]{article}
\usepackage[margin=1in]{geometry}
\usepackage{amsmath,amssymb,amsthm}
\usepackage{booktabs}
\usepackage{tabularx}
\usepackage[pdftitle={Supplementary Material: Corrected Statements and Proofs
for the Revised Long-Memory Limit Theorems}, pdfauthor={hurst_function_
estimation Lean audit}]{hyperref}

\newtheorem{theorem}{Theorem}
\newtheorem{proposition}[theorem]{Proposition}
\newtheorem{lemma}[theorem]{Lemma}
\theoremstyle{definition}
\newtheorem{corollary}[theorem]{Corollary}
\newtheorem{condition}[theorem]{Condition}
\theoremstyle{remark}
\newtheorem{remark}[theorem]{Remark}

\newcommand{\statustag}[1]{%
  \par\nobreak\begingroup\footnotesize\itshape\raggedleft #1\par\endgroup}
\newcommand{\Wlean}{[Lean]}
\newcommand{\Wwritten}{[written]}

\title{Supplementary Material\\
\large Corrected Statements and Proofs for the Revised Long-Memory
Limit Theorems\\
\footnotesize Companion to \texttt{revised\_theorems.tex}; machine-checked
items verified in the Lean~4 repository at closure commit \texttt{16715db}}
\author{}
\date{October 7, 2026}

\begin{document}
\maketitle

\noindent This supplement gives, for each substantive correction identified in
the audit of \texttt{19-AOS1825.pdf}: the original claim, the error (with the
counterexample where one exists), the corrected statement, and a proof or
proof sketch. Section~\ref{sec:drift} then proves in full the honest-rate
drift proposition underlying the revised main theorem. Full written chains are
in \texttt{direct\_proofs}/01--24; per-item machine status is in
\texttt{paper\_revision/theorem\_inventory.md}.
Status tags: \Wlean{} = completely formalized;
\Wwritten{} = written mathematical proof. Proofs marked ``English rendering''
are faithful translations of the corresponding Chinese sources in
\texttt{direct\_proofs}/; every item also carries a pointer
\emph{Full written chain} to the exact source files, which remain authoritative
for any detail not reproduced here.

\section{Corrections to the main theorems}\label{sec:main-corr}

\subsection{Critical-regime variance (Theorem 3.3(ii), via S.3.2)}

\noindent\textbf{Original.} In the critical regime $H(t)=\tfrac34$ the limiting
variance is expressed through the kernel center value.
\textbf{Error.} The derivation of the diagonal contribution $\varepsilon_4$
in S.3.2 replaces $f^2(i/N)f^2(j/N)$ by $f^4(0)$ for $|i-j|$ small, which is
not summable to $f^4(0)$; the correct diagonal sum tends to $\int f^2$, and
the Hermite rank-$2$ factor $2!$ together with the squared denominator of
$g$ was mishandled.
\textbf{Corrected.}
\[
\operatorname{Var} \;\longrightarrow\; \frac{9}{16}\int \omega_r(x)^2\,dx .
\]
\textbf{Proof (English rendering of \texttt{direct\_proofs}/03, \S5--6).}
For unit-stride first increments of fractional Brownian motion with constant
$H$,
\[
r_H(k)=\tfrac12\bigl(|k+1|^{2H}+|k-1|^{2H}-2|k|^{2H}\bigr),\qquad r_H(0)=1 .
\]
At $H=\tfrac34$ a Taylor expansion gives
$r_{3/4}(k)=\tfrac38|k|^{-1/2}+O(|k|^{-5/2})$, so both long-memory spectral
coefficients equal $a(\pm1)=\tfrac38$ and $c_a=2(\tfrac38)^2=\tfrac9{32}$.
The log-square statistic is a weighted quadratic form in the standardized
increments; diagonalizing it as $Q_N=\sum_j\lambda_{N,j}(Z_j^2-1)$ with
$s_N^2=\sum_j\lambda_{N,j}^2$, the operator bound
$\max_j|\lambda_{N,j}|\le\|B_N\|_{\mathrm{op}}=O(N^{1/2})$ gives
$\max_j|\lambda_{N,j}|/s_N=O((\log N)^{-1/2})\to0$, and the
independent-$\chi^2$ log-characteristic function
\[
\log \mathbb E e^{itQ_N/s_N}
=\sum_j\Bigl\{-itu_{N,j}-\tfrac12\log(1-2itu_{N,j})\Bigr\}
=-\tfrac{t^2}{2}+o(1)
\]
(with remainder $\le C_t\max_j|u_{N,j}|\sum_ju_{N,j}^2\to0$) yields the
Gaussian limit by L\'evy continuity --- notably \emph{without} invoking the
short-memory Breuer--Major theorem, whose summability hypothesis fails at
criticality. The limiting variance is $2c_a\int f^2=\frac9{16}\int\omega^2$.

\smallskip
\emph{Counterexample to the printed formula.} Take $d=q=p=1$, $H\equiv
\frac34$, and the smooth nonnegative bump
$K_0(z)=z^2e^{-1/(1-z^2)}$ on $|z|<1$ (zero otherwise), normalized to
$K=K_0/\int K_0$, with $t$ interior and $b=n^{-1/2}$. The equivalent kernel
$\omega=K/\int K$ then satisfies $\omega(0)=0$ but $\int\omega^2>0$. By fBm
stationarity and self-similarity the standardized local increment vectors
carry exactly the covariance $r_{3/4}$; the actual weights differ from
$N^{-1}\omega(z_i)$ by $N^{-1}e_{N,i}$ with $\sup|e_{N,i}|=O(N^{-1})$, and
the covariance bound $|\operatorname{Cov}(g(Y_i),g(Y_j))|\le Cr(i-j)^2$ makes
the weight-perturbation variance $O(N^{-2})$ after critical normalization.
Hence the centered $\widehat G$ converges under the paper's
$\sqrt{N/\log N}$ normalization to a normal law with variance
$\frac9{16}\int\omega^2>0$, while the paper's kernel-center-value formula
returns $0$: a counterexample to the \emph{universal critical variance
formula itself}, not merely to a step of its derivation.
\statustag{\Wwritten}

\smallskip\noindent{\footnotesize Full written chain:
\texttt{direct\_proofs/03}, \texttt{direct\_proofs/13} (Chinese sources).}

\subsection{Long-memory sign (Theorem 3.3(iii)--3.4)}\label{sec:sign}

\noindent\textbf{Original.} The inverted estimator converges to $+Q$.
\textbf{Error.} The inversion $x\mapsto(c_\sigma-x)/(2\log n)$ carries a
negative slope, so the limit reflects: $\,-Q$. This is substantive: a legal
nonnegative equivalent-kernel instance has third moments of $Q$ and $-Q$ of
opposite sign, so the two laws differ.
\textbf{Corrected.}
$2S_n^{\psi}\log n\,(\widehat H_n-h)\Rightarrow -Q$
(Theorem~1 of the companion note).
\textbf{Proof (direct\_proofs/14, \S7; the full statement is machine-checked).}
The calibration $x\mapsto(c_\sigma-x)/(2\log n)$ is decreasing, so the
expectation-centered limit of $\widehat G$ transports to the \emph{reflected}
limit $-Q$; the affine inversion argument at $a_n=N^{\psi}$ carries the
deterministic center to the true expectation center using
$V_n=O(N^{-2\psi})$.

The sign is substantive. Take $p=2$ and a nonnegative, symmetric, nonzero
kernel $K$; the local-linear equivalent kernel $\omega=K/\int K$ is then
nonnegative, and for any admissible $H$ with $H(t)>\frac34$ (constant $H$
suffices) the limiting quadratic-form operator $\mathcal A$ is nonnegative
and nonzero: using the representation
$\tfrac12e^{-a|x-y|}=a\int_{-\infty}^{\min(x,y)}e^{-a(x-u)}e^{-a(y-u)}du$,
the $\omega$-weighted quadratic integral is an integral of nonnegative
squares, and if it vanished then
$F(u)=\int_u^\infty\omega(x)e^{-a(x-u)}dx\equiv0$ together with
$F'=aF-\omega$ would force $\omega=0$, contradicting $\int\omega=1$
(all exchanges are justified by the integrability of the kernel family).
Hence $\operatorname{tr}(\mathcal A^3)>0$ and
\[
\mathbb EQ^3=8\operatorname{tr}(\mathcal A^3)>0
\qquad\text{while}\qquad
\mathbb E(-Q)^3<0,
\]
so $Q$ and $-Q$ are not identically distributed --- on a submodel allowed by
the paper's assumptions, at the level of the actual statistic's limit. The
fourth cumulant $48\sum_j\lambda_j^4>0$ ($\mathcal A\neq0$) also shows $Q$ is
not Gaussian.
\statustag{\Wlean}

\smallskip\noindent{\footnotesize Full written chain:
\texttt{direct\_proofs/14}, \texttt{direct\_proofs/13} (Chinese sources).}

\subsection{Optimal-bandwidth biased means (Corollary 3.5)}
\textbf{Proof (English rendering of \texttt{direct\_proofs}/13, \S5--6).}
Write $\mu_n=\mathbb E\widehat G(t)$, $\theta_n=(c_\sigma-\mu_n)/(2L)$ with
$L=\log n$. The file-12 variance bound gives the true variance
$V_n=O(a_n^{-2})$ and $\theta_n\to h$. On the event that the truncation does
not act, the \emph{exact} identity
\[
2a_nL(\widehat H-\theta_n)=-a_n(\widehat G-\mu_n)
\]
holds; the truncation event is contained in
$\{|\widehat G-\mu_n|\ge cL\}$, of probability at most $CV_n/L^2\to0$, and
$a_nL\,|\mathbb E\widehat H-\theta_n|\le Ca_nV_n/L=O((a_nL)^{-1})\to0$,
which carries the deterministic center to the true expectation center and
keeps the sign.

For the biased means, let $p$ be an integer and let $H$ have a $p$-th order
Peano expansion at $t$, and set
$R_H(t)=\frac{H^{(p)}(t)}{p!}\int z^p\omega(z)\,dz$. Weight reproduction
together with file 12 gives
\[
B_n:=\mathbb E\widehat G-G_n(h)=-2Lb^{p}R_H(t)+o(Lb^{p})+O(\rho_n(t)) .
\]
In the short-memory regime $h<\tfrac34$ take $b_0=(nL^2)^{-1/(2p+1)}$
with $b/b_0\to c_0\in(0,\infty)$, and set $A_n=(n^{p}/L)^{1/(2p+1)}$; then
$A_nLb_0^{p}=1$ and $A_n\rho_n(t)\to0$ because
$\psi>\tfrac12>p/(2p+1)$. Consequently
\[
A_n(\widehat G-G_n(h))\Rightarrow
N\bigl(-2c_0^{p}R_H(t),\,c_0^{-1}V_h\bigr),
\qquad
2(nL^2)^{p/(2p+1)}(\widehat H-h)\Rightarrow
N\bigl(2c_0^{p}R_H(t),\,c_0^{-1}V_h\bigr),
\]
the sign flip coming from the decreasing calibration. When $R=0$ the
additive $o$-term from the bias expansion is used; no equivalence with $0$
is claimed.
\statustag{\Wwritten; conditional Lean chain connected}

\smallskip\noindent{\footnotesize Full written chain:
\texttt{direct\_proofs/13}, \texttt{direct\_proofs/11},
\texttt{direct\_proofs/01} (Chinese sources).}

\subsection{Growing-lag remainder (Theorem 8.2(iv))}
\textbf{Counterexample (English rendering of \texttt{direct\_proofs}/04,
\S2).} Take constant $H=\tfrac34$, $q=1$, $\sigma=1$. By stationarity and
self-similarity of fBm increments, the paper's standardized difference
covariance is exactly $g_1(\tfrac34,u)$, with no $n$-dependent freezing
error. A sixth-order Taylor expansion of $(u\pm1)^{3/2}$ (odd terms
cancelling; coefficients $\binom{3/2}{2}=\tfrac38$,
$\binom{3/2}{4}=\tfrac3{128}$; remainder controlled by derivative bounds for
$u\ge2$) gives
\[
g_1(\tfrac34,u)\,|u|^{1/2}=\frac38+\frac3{128u^{2}}+O(u^{-4}),
\]
so for all $u\ge u_0$ the left side exceeds $\tfrac38$ by at least
$\tfrac3{256u^2}$. Now set $b_n=n^{-1/2}$ and $u_n=\lfloor\log n\rfloor$ at a
fixed interior $t$: then $nb_n\to\infty$, $b_n\log^{k}n\to0$ for every fixed
$k$, $u_n\to\infty$, $u_n<2nb_n$, and $t+u_n/n$ stays interior --- all the
hypotheses of the original 8.2(iv) hold. But
\[
\frac{u_n^{-2}}{b_n|\log b_n|}\asymp\frac{\sqrt n}{(\log n)^{3}}\to\infty
\qquad(n=e^{x}\text{ swamps the cubic}),
\]
so the remainder cannot be $O(b\log b)$: the original 8.2(iv) fails already
on the constant-$H$ submodel allowed by the paper.

\textbf{Corrected versions.} (i) Constant $H$:
$C_n(t,t+u/n)=\{c_q(H)+O(|u|^{-2})\}|u|^{-\psi}$, $\psi=2q-2H$.
(ii) Nonconstant $H$ (English rendering of \texttt{direct\_proofs}/12,
\S3.3): for $2<|k|\le2nb$, from the integral representation, the H\"older
increment $H(s+v)-H(s)=O(1/n)$, and the normalization factor
$n^{H(u)-H(s)}=1+O(|k|L/n)$,
\[
C_n(s,s+k/n)=\Bigl\{c_1(H(s))+O\Bigl(k^{-2}+\tfrac{|k|L}{n}
+(|k|/n)^{\psi(s)}\Bigr)\Bigr\}|k|^{-\psi(s)} .
\]
The mechanism: the power error from the shifted exponent is $O(L/n)$, the
derivative distance-logarithm remainder at most $C|k|L/n$, the smooth
background contributes $(|k|/n)^{\psi(s)}$, and the frozen integral, written
with $V=n(w-v)$ satisfying $\mathbb EV=0$, $|V|\le1$, obeys the second-order
Taylor bound
$\mathbb E|k+V|^{-\psi(s)}=|k|^{-\psi(s)}(1+O(k^{-2}))$. All remainders are
absolute coefficient errors, never divided by the possibly-vanishing $c_1$.
\statustag{\Wwritten; interior pieces Lean}

\smallskip\noindent{\footnotesize Full written chain:
\texttt{direct\_proofs/04}, \texttt{direct\_proofs/12},
\texttt{direct\_proofs/18} (Chinese sources).}

\subsection{Local-polynomial bias (Lemma 8.4)}\label{sec:84}

\noindent\textbf{Original.} An integral-region/quadrature argument.
\textbf{Error.} The quadrature over the integration region is invalid.
\textbf{Corrected.} Define $M_{kl}=\sum_i K_i A_{ik}A_{il}$ and
$w_i = K_i\sum_l (M^{-1})_{0l}A_{il}$.
\textbf{Proof.} Invertibility of $M$ gives
$\sum_i w_i A_{ik}=\delta_{k0}$ exactly; hence for
$F_i=\beta_0+\sum_{k\ge1}\beta_k A_{ik}+r_i$ with $|r_i|\le B$ on the support
of $K$,
\[
\Bigl|\sum_i w_iF_i-\beta_0\Bigr|\le \sum_i|w_i|\,B = O(b^p),
\]
using the uniform $L^1$ weight bound. No $(nb)^{-\min(2,p)}$ term is needed.
\statustag{\Wlean (designs, weights, equivalent kernel, bias leading term)}

\smallskip\noindent{\footnotesize Full written chain:
\texttt{direct\_proofs/02}, \texttt{direct\_proofs/18} (Chinese sources).}

\subsection{Vector $Z$ limits (Lemma 8.5)}
\textbf{Corrections.} (i) The three limit statements must be made for
$Z-\mathbb EZ$: the normalization of the estimator chain transports centered
vectors, and the bound $\mathbb EZ=O((nb)^{d/2}\rho)$ alone does not prove
that the normalized $\mathbb EZ$ vanishes --- if the original uncentered $Z$
is kept, this vanishing must be established separately.
(ii) The dimension of the local-polynomial coefficient vector $a$ is the
basis dimension $S=\binom{p+d}{d}$-type count, not the ambient $d$; the
covariance structure and the limits are stated in $\mathbb R^{S}$.
(iii) Kernel conditions ensuring the nondegeneracy used by the limits are
added explicitly (they are exactly the ones repaired in S.2.1 below).
The three limit statements themselves are at the written level (the
finite-Hermite CLT for non-absolutely-summable correlation arrays of file 13
\S2 supplies the mechanism); the expectation-bias leading terms derived from
them are machine-checked.
\statustag{expectation-bias leading terms \Wlean; three $Z$-limits \Wwritten}

\smallskip\noindent{\footnotesize Full written chain:
\texttt{direct\_proofs/03}, \texttt{direct\_proofs/13},
\texttt{direct\_proofs/14} (Chinese sources).}

\section{Corrections to the supplementary lemmas}\label{sec:supp-corr}

\subsection{S.1.1 (Fourier adjacent differences)}
\textbf{Diagnosis.} As a complex Lebesgue integral the claim fails: near
$0$ the imaginary part of the integrand behaves like $-ikx/|x|^{2+\delta}$,
which is not absolutely integrable.

\textbf{Corrected statement and proof (English rendering of
\texttt{direct\_proofs}/06, \S2).} For $0\le\delta\le\tfrac12$ define
\[
A_\delta=\int_{\mathbb R}\frac{1-\cos x}{|x|^{2+\delta}}dx,
\qquad
B_\delta=\int_{\mathbb R}\frac{(1-\cos x)\log|x|}{|x|^{2+\delta}}dx .
\]
$A_\delta$, and the absolute integral defining $B_\delta$, are uniformly
bounded on this $\delta$-range: near $0$ they are dominated by
$C|x|^{-1/2}$ and $C|x|^{-1/2}|\log|x||$, in the tail by $C|x|^{-2}$ and
$C|x|^{-2}\log|x|$. For $k\ge1$ the substitution $y=kx$ gives the exact
scalings
\[
I_\delta(k)=\int\frac{1-\cos(kx)}{|x|^{2+\delta}}dx=A_\delta\,k^{1+\delta},
\qquad
J_\delta(k)=\int\frac{(1-\cos(kx))\log|x|}{|x|^{2+\delta}}dx
=k^{1+\delta}\bigl(B_\delta-A_\delta\log k\bigr),
\]
with $I_\delta(0)=J_\delta(0)=0$. Assume $\delta_n\ge0$ and
$n^{\delta_n}\to1$; then eventually $\delta_n\le\tfrac12$ and
$n^{\delta_n}\le2$. For $2\le k\le n$ the mean-value theorem gives
$k^{1+\delta_n}-(k-1)^{1+\delta_n}\le(1+\delta_n)n^{\delta_n}\le3$, and the
derivative of $x^{1+\delta}\log x$ is
$x^{\delta}[(1+\delta)\log x+1]=O(\log n)$ on $[1,n]$. With the uniform
bounds on $A,B$,
\[
\sup_{1\le k\le n}|I_{\delta_n}(k)-I_{\delta_n}(k-1)|=O(1),
\qquad
\sup_{1\le k\le n}|J_{\delta_n}(k)-J_{\delta_n}(k-1)|=O(\log n),
\]
the case $k=1$ being handled separately from the boundedness of $A_\delta$,
$B_\delta$. This proves the real-part (cosine) repaired version of S.1.1;
under symmetric truncation the imaginary part vanishes identically, so the
same proof covers the explicitly defined principal-value version.
\statustag{\Wwritten (repair); not formalized}

\smallskip\noindent{\footnotesize Full written chain:
\texttt{direct\_proofs/06} (Chinese source).}

\subsection{S.1.2 (filtered Fourier integral)}
\textbf{Status.} No counterexample found; a more direct route exists but is
not formalized. \textbf{Route.} Use the cosine form and the second difference
$(k+1)^{1+\delta}-2k^{1+\delta}+(k-1)^{1+\delta}$ in place of the contour
argument; $k=1$ is handled separately and $\delta=0$ returns $0$; the real
part must be retained in the contour derivation.

\smallskip\noindent{\footnotesize Full written chain:
\texttt{direct\_proofs/06} (Chinese source).}

\subsection{S.2.1 ($g$ bounded away from zero)}
\textbf{Error.} The hypotheses are insufficient: at $h=0$ the difference
vanishes identically, and without compactness away from $1$ there is no
uniform positive lower bound --- for $q=2$,
$g_2(H,0,1)=4-2^{2H}\to0$ as $H\uparrow1$.

\textbf{Corrected statement and proof (English rendering of
\texttt{direct\_proofs}/05, \S2).} Fix $q\ge1$ and $h\neq0$. For the
constant-$H$ model, applying $q$-th order differences on both sides of the
corrected spectral representation (5.2) of \S S.2.2 below gives
\[
g_q(H,0,h)=\frac{1}{2D(H)^{2}}
\int_{\mathbb R^{d}}\frac{|1-e^{i\langle h,\xi\rangle}|^{2q}}{|\xi|^{2H+d}}
\,d\xi\;>\;0 :
\]
near $0$ the numerator is $O(|\xi|^{2q})$ and the integral converges because
$q>H$; the tail converges because $H>0$; and $h\neq0$ makes the numerator
strictly positive on an open set. If $H\in[\eta,1-\eta]$, the same
integrable majorant controls the whole family (radially
$r^{2q-2(1-\eta)-1}$ near $0$ and $r^{-2\eta-1}$ at infinity), and similarly
for $D(H)$; dominated convergence makes $H\mapsto g_q(H,0,h)$ continuous,
whence by compactness
$\inf_{H\in[\eta,1-\eta]}g_q(H,0,h)>0$. This is a direct repair that does
not rely on the conditionally-negative-definite kernel theorem quoted in the
original; the two added conditions ($h\neq0$; compactness away from $0$ and
$1$, available from [A1] rather than [A3]) are both indispensable, as the
two counterexamples show.
\statustag{repaired 1-d version \Wlean}

\smallskip\noindent{\footnotesize Full written chain:
\texttt{direct\_proofs/05} (Chinese source).}

\subsection{S.2.2 / S.3.1 (leading singular coefficient)}
\textbf{Proof of the corrected constants} (English rendering of
\texttt{direct\_\allowbreak{}proofs}/\allowbreak{}04, \S1, and of
\texttt{direct\_\allowbreak{}proofs}/\allowbreak{}05, \S1).
(i) \emph{Difference expansion.} Let $q\ge1$ and $F_H(u)=|u|^{2H}$, and
\[
g_q(H,u)=-\tfrac12\sum_{i,j}(-1)^{i+j}\binom qi\binom qj
F_H(u+i-j).
\]
By the binomial expansion (Vandermonde),
$g_q(H,u)=\tfrac{(-1)^{q+1}}2\,\delta^{2q}F_H(u)$ with
$\delta F(u)=F(u+\tfrac12)-F(u-\tfrac12)$. For $u>2q$, repeated integration
of the fundamental theorem gives
$\delta^{2q}F_H(u)=\int_{[-1/2,1/2]^{2q}}F_H^{(2q)}(u+\textstyle\sum s_i)\,ds$;
the mean of $S=\sum s_i$ is $0$ with $|S|\le q$, and a first-order Taylor
expansion of $F_H^{(2q)}$ has second-order remainder bounded by
$Cu^{2H-2q-2}$ uniformly for $H$ in a fixed compact subinterval. Hence
\[
g_q(H,u)=c_q(H)|u|^{2H-2q}+O(|u|^{2H-2q-2}),
\qquad
c_q(H)=\frac{(-1)^{q+1}}{2}\prod_{l=0}^{2q-1}(2H-l),
\]
and in particular $c_1(H)=H(2H-1)$: the paper's printed $c_h$ is missing
the factor $\tfrac12$. When the product vanishes, (the display) remains an
absolute remainder bound --- nothing is divided by $c_q(H)$. The far-field
form $g=(c_h+o(1))|u|^{-\psi}$ of S.3.1 follows with the same $c_h$, and
its $o(1)$ must not be strengthened to $O(b\log b)$.
(ii) \emph{Spectral normalization.} With
$D(H)^2=\int(1-\cos\xi_1)/|\xi|^{2H+d}\,d\xi$ (finite and positive for
$0<H<1$), rotation and scaling give
$\int(1-\cos\langle t,\xi\rangle)/|\xi|^{2H+d}d\xi=D(H)^2|t|^{2H}$. For
$a=\frac{H_t+H_s}{2}$ the absolutely convergent integral
$I(t,s)=\int(e^{i\langle t,\xi\rangle}-1)(e^{-i\langle s,\xi\rangle}-1)
/|\xi|^{H_t+H_s+d}\,d\xi$ has vanishing imaginary part (odd symmetry) and
real part $D(a)^2\{|t|^{2a}+|s|^{2a}-|t-s|^{2a}\}$, so the correct
covariance normalization is
$C(t,s)=\frac{\sigma^2}{2D(H_t)D(H_s)}I(t,s)$: omitting the $\tfrac12$
would double the variance to $2\sigma^2|t|^{2H}$ at $t=s$.
\statustag{instantiated and machine-verified inside the closure chain \Wlean}

\smallskip\noindent{\footnotesize Full written chain:
\texttt{direct\_proofs/04}, \texttt{direct\_proofs/05},
\texttt{direct\_proofs/12} (Chinese sources).}

\subsection{S.2.3 (mixed derivatives, near-diagonal)}
\textbf{Counterexample (English rendering of \texttt{direct\_proofs}/12,
\S2).} Fix an admissible base point $(h_0,t_0)$ with $t_0$ interior and
$h_0\in(0,1)$, and take the lawful smooth Hurst function
$H_\varepsilon(x)=h_0+\varepsilon(x-t_0)$ with $\varepsilon>0$ small (it is
smooth and stays away from $0$ and $1$). Let the smooth background be
$P(s,u)=B(H_\varepsilon(s),H_\varepsilon(u))(s^{a}+u^{a})$, where $B$ is the
normalizing factor of the representation. $B$ is symmetric with
$B(h,h)=\tfrac12$, hence $B_h(h,h)=B_k(h,h)=0$, and direct differentiation
gives
\[
\partial_s\partial_uP(t_0,t_0)
=\varepsilon t_0^{2h_0-1}\bigl(2h_0\log t_0+1\bigr)
+\varepsilon^{2}t_0^{2h_0}\bigl\{2B_{hk}(h_0,h_0)+(\log t_0)^{2}\bigr\}
=\tfrac12\varepsilon t_0^{2h_0-1}+O(\varepsilon^{2})\;>\;0 .
\]
Fix such an $\varepsilon$. The remainder of the singular part alone,
computed in \S2 of the file, is $O(r^{1-\psi_0}(1+|\log r|))=o(1)$ with
$\psi_0=\tfrac12$, so
\[
\lim_{r\downarrow0}\Bigl\{\partial_s\partial_uC(t_0,t_0+r)
-c_1(h_0)\,r^{-\psi_0}\Bigr\}=\partial_s\partial_uP(t_0,t_0)>0 ,
\]
whereas the original S.2.3 remainder, a \emph{relative} $O(r\log r)$, would
force the bracketed difference to tend to $0$ --- a contradiction. The
principal-term expansion as printed therefore fails.

\textbf{Repair.} Add the smooth background of relative size $r^{\psi}$: the
corrected expansion retains, near the diagonal, the
$x^{H(x)+H(y)-q}$-times-log-power background derivative, and the normalized
finite differences on the actual grid are then controlled by it; when
$H(s)>\tfrac12$ one has $\psi(s)<1$ and the term cannot be absorbed into
$O(r|\log r|)$. Files 12/13/14 verify that this repair does not change the
one-dimensional variance rates or limit types.
\statustag{\Wwritten; grid normalization pieces \Wlean}

\smallskip\noindent{\footnotesize Full written chain:
\texttt{direct\_proofs/12}, \texttt{direct\_proofs/18},
\texttt{direct\_proofs/04} (Chinese sources).}

\subsection{S.3.2 (weighted power sums / cycle products)}
Corrected jointly with \S1.1 above (center value $\int f^2$; $2!$;
denominator). The long-memory side --- identification of the cycle products
with the spectrum of $B=K^{1/2}M_\omega K^{1/2}$ for all orders $k\ge2$ ---
is machine-checked. \statustag{long-memory identification \Wlean}

\smallskip\noindent{\footnotesize Full written chain:
\texttt{direct\_proofs/03}, \texttt{direct\_proofs/13},
\texttt{direct\_proofs/14} (Chinese sources); the critical side is proved in
\S1.1 above.}

\subsection{S.5.1 (spatial average variance of $I_1'$)}
\textbf{Error.} The far-distance count misses a factor $m$: with
$H\equiv7/8$, $q=p=1$, $b=n^{-1/2}$ one has
$\operatorname{Var}(I_1)\approx\log^2n\cdot n^{-1/2}$ against the claimed
$O(\log^2n\cdot n^{-3/4})$.
\textbf{Correction.} The count is $m^{2\psi-1}n^{-2\psi}$ times the lag sum;
merging weights at $m\approx1/b$ gives
$\operatorname{Var}(I_1)\le C\log^2n\cdot T(n)$, which the original bandwidth
condition still absorbs into the target term --- the main scale rates are
preserved. \statustag{\Wwritten; the $q=1$ scale input it feeds is
discharged in Lean}

\subsection{S.5.2 (truncated component $I_2'$)}
\textbf{Proof of the repaired bounds (English rendering of
\texttt{direct\_proofs}/16, \S4).} The pilot supplies, uniformly in $j$,
\[
|\mathbb EP(s_j)-H(s_j)|\le Ca_n,
\qquad
\mathbb E|P(s_j)-H(s_j)|^{2}\le C\{a_n^{2}+T_{\bar\psi}(N)\}.
\]
For $q=1$ the truncation functional $\ell$ vanishes identically, so
$I_2=0$. For $q=2$, $\ell,\ell',\ell''$ are bounded on
$[0,1-\gamma/2]$ and the true $H$ keeps a fixed positive distance from both
truncation endpoints. Put $v=\mathbb E(P-H)^{2}$. Projection does not
increase the error relative to $H$, so
$\mathbb E|\ell(P^{\gamma})-\ell(H)|^{2}\le Cv$; moreover
$|P^{\gamma}-P|\le|P-H|\mathbf 1_{\{|P-H|\ge d_0\}}$ gives
$\mathbb E|P^{\gamma}-P|\le v/d_0$; and Taylor's formula gives
$|\mathbb E\{\ell(P^{\gamma})-\ell(H)\}|\le C\{|\mathbb E(P-H)|+v\}$.
Applying Jensen to the spatial average --- with \emph{no} independence
assumed across $j$, hence no fictitious $1/M$ variance gain --- yields
\[
\mathbb EI_2^{2}\le C\{a_n^{2}+T_{\bar\psi}(N)\},
\qquad
|\mathbb EI_2|\le C\{a_n+T_{\bar\psi}(N)\},
\]
and $a_n^{2}$ is absorbed by $a_n$ since $a_n\to0$. Combined with the
repaired S.5.1 bound and the square-sum inequality this recovers Theorem
4.2's MSE decomposition; the role of the bandwidth condition
$L^{2}T(n)=o(T(nb_1))$ is checked separately in each memory regime
(short: $bL^{2}\to0$; long: $L^{2}b^{2\bar\psi}\to0$ by choosing the fixed
$k$ of (16.1) large; critical: at most $2bL^{2}$ if $N\ge\sqrt n$, at most
$n^{-1/2}L^{3}/\log N$ otherwise).
\statustag{\Wwritten + Full pieces}

\smallskip\noindent{\footnotesize Full written chain:
\texttt{direct\_proofs/16}, \texttt{direct\_proofs/08} (Chinese sources).}

\subsection{S.6.1 (irregular-grid covariance)}
\textbf{Counterexample (English rendering of \texttt{direct\_proofs}/05,
\S4).} Take one-dimensional standard Brownian motion $B$, $q=1$,
$H=\tfrac12$; regular grids are admissible special cases of the irregular
grids allowed. By the Brownian covariance $\min(s,t)$,
$\operatorname{Var}\{B(t+1/n)-B(t)\}=1/n$, while
$g_1(\tfrac12,0,1)=1$; the difference of the two sides tends to $-1$ and
cannot be $O(\log n/n+1/n)$. This refutes the unnormalized claim S.6.1(i)
as printed. (The scalar asymptotic underlying this counterexample is
machine-checked in the repository.)

\textbf{Correct transformation.} For arbitrary grid differences
$D_n(t)$ define $W_n(t)=\sigma^{-1}n^{H(t)}D_n(t)$. Bilinearity of the
covariance gives the \emph{exact} identity
\[
\operatorname{Cov}\bigl(D_n(t),D_n(s)\bigr)
=\sigma^{2}n^{-H(t)-H(s)}
\operatorname{Cov}\bigl(W_n(t),W_n(s)\bigr),
\]
so if the normalized version states
$\operatorname{Cov}(W_n(t),W_n(s))=g(H(t),u,\varpi(t))+O(\rho_n)$, the
original-difference version must be
$\sigma^{2}n^{-H(t)-H(s)}\{g(H(t),u,\varpi(t))+O(\rho_n)\}$ --- the error
term carries the same prefactor. On the Brownian regular grid with $u=0$
the corrected form is an exact equality; for general irregular grids the
local covariance approximation still requires separate proof --- the
identity alone does not provide it.
\statustag{scalar counterexample \Wlean; full instantiation open}

\smallskip\noindent{\footnotesize Full written chain:
\texttt{direct\_proofs/05} (Chinese source).}

\subsection{S.7.2 (dimension scope)}
No counterexample; scope remark: the $d=2,3$ ``minimax optimality'' assertion
has no corresponding lower bound in the manuscript itself, so it cannot be
certified even after repairing the upper-bound route.
\statustag{\Wwritten (route only)}

\smallskip\noindent{\footnotesize Full written chain:
\texttt{direct\_proofs/02} (Chinese source).}

\section{Proof of the honest-rate drift proposition}\label{sec:drift}

This section proves Proposition~2 of the companion note in full. Work under
(R1)--(R6) of the companion note: $\delta_n=n^{-h}$, $S_n=n^{1-h}$,
$\psi=2-2h\in(0,\tfrac12)$, $c_\sigma=\mathbb E\log Z^2$, $\widehat G_n$ the
unit-weight log-square statistic, $\widetilde H_n=(c_\sigma-\widehat
G_n)/(2\log n)$, $\widehat H_n=T(\widetilde H_n)$ clipped to $[0,1]$.

\begin{proposition}\label{prop:drift-full}
$2S_n^{\psi}\log n\,\bigl|\mathbb E\widehat H_n-h\bigr|\to0$, with the
three-term bound of the companion note; the second term vanishes exactly
under the $b$-band $(1-h)^2<1-b$, the other two unconditionally.
\end{proposition}

\begin{proof}
\textbf{Step 1 (exact expectation identity).} For each active row $j$ the
coefficient-feature combination satisfies, by the stride identity, the
norm decomposition
\[
\log\|\text{combo}_j\|^2=-2H_j\log n+\log\|V_j\|^2,
\]
where $H_j=f(s_j)$ is the left-endpoint H\"older value and $V_j$ the
normalized varying increment. The Gaussian expectation formula gives
\[
\mathbb E\widehat G_n=\sum_j w_j\bigl(-2H_j\log n+\log\|V_j\|^2\bigr)+c_\sigma,
\qquad \sum_j w_j=1,
\]
whence, with $c_\sigma$ pinned,
\begin{equation}
\mathbb E\widetilde H_n-h
= -\Bigl(\sum_j w_jH_j-h\Bigr)-\frac{1}{2\log n}\sum_j w_j\log\|V_j\|^2 .
\label{eq:decomp}
\end{equation}

\textbf{Step 2 (deterministic side).}
(i) \emph{Weight contraction.} On the active window $|s_j-t|<\delta_n$, so
H\"older regularity gives $|H_j-h|\le D(1+M)\delta_n$, and
$\bigl|\sum_jw_jH_j-h\bigr|\le\sum_j|w_j||H_j-h|\le C_{\mathrm{st}}\delta_n^{p}$
with $p\ge1$ (the sharper exponent uses the class regularity; $p=1$ suffices).
(ii) \emph{Second-order increment log rate.} Write $V_j=F_{k}(s_j+\ell)
$-type normalized varying increment with $|k-h|\le K\ell$, $\ell=1/n$,
$K=D(1+M)$. The exact kernel identity
\[
2\langle \text{frozen},\,\text{mismatch}\rangle
=\Lambda\,\Phi(m)-\Phi(h),\qquad
\Phi(q)=(s+\ell)^{2q}-s^{2q}+\ell^{2q},\quad m=\tfrac{h+k}{2},
\]
with $\Lambda-1=-\tfrac12\|F_k(1)-F_h(1)\|^2=O((k-h)^2)$ (uniform parameter
Lipschitz; no smoothness of the normalizing constant needed) and the
mixed-power Lipschitz bound $|\Phi(m)-\Phi(h)|\le C|k-h|\,\ell$ (valid since
$2a>1$ absorbs the logarithmic sensitivity) yields
\[
\bigl|\|V_j\|^2-1\bigr|\le C(1+K+K^2+K^3)\,\ell^{\,2-2b}
=C\,n^{-(2-2b)},
\]
hence $|\log\|V_j\|^2|\le 2C\,n^{-(2-2b)}=:C_{\log}n^{-(2-2b)}$ once small,
and $|\sum_jw_j\log\|V_j\|^2|\le C_wC_{\log}n^{-(2-2b)}$
with $C_w=\sum_j|w_j|$ uniformly bounded.

\textbf{Step 3 (random side; tail-based transfer).} Let
$m=\mathbb E\widetilde H_n$. Since $\mathbb E\widetilde H_n\to h\in(3/4,1)$
(companion note, center-band discharge), eventually
$m\in[d',1-d']$ with $d'=(1-h)/2$. On $\{\widetilde H_n\notin[0,1]\}$ we have
$|\widetilde H_n-m|\ge d'$. The unit-weight variance bound and the normalized
energy give
\[
V_n:=\operatorname{Var}(\widetilde H_n)
=\frac{\operatorname{Var}(\widehat G_n)}{4\log^2n}
\le\frac{4\,g_{\mathrm{LSV}}\,U^2C_2}{4\log^2n}\,S_n^{-2\psi}
=O\!\left(\frac{n^{-4(1-h)^2}}{\log^2n}\right),
\]
and Chebyshev gives $\varepsilon_n:=\mathbb P(|\widetilde H_n-m|\ge d')\le
V_n/d'^2$. Since $|T(x)-x|\le|x-m|$ off $[0,1]$ and $T(x)=x$ on $[0,1]$,
\[
\bigl|\mathbb E\widehat H_n-\mathbb E\widetilde H_n\bigr|
\le\mathbb E\bigl[|\widetilde H_n-m|\mathbf 1_{\widetilde H_n\notin[0,1]}\bigr]
\le V_n+\varepsilon_n \le V_n\bigl(1+d'^{-2}\bigr).
\]

\textbf{Step 4 (exponents).} Combining
$\mathbb E\widehat H_n-h=[\mathbb E\widetilde H_n-h]+[\mathbb E\widehat
H_n-\mathbb E\widetilde H_n]$ with Steps 2--3,
\[
2S_n^{\psi}\log n\,\bigl|\mathbb E\widehat H_n-h\bigr|
\le
\underbrace{C_1\,\frac{\log n}{n^{\,h-2(1-h)^2}}}_{\text{smoothing}}
+\underbrace{C_2\,n^{-(2-2b)-2(1-h)^2}}_{\text{mismatch}}
+\underbrace{C_3\,\frac{1}{n^{2(1-h)^2}\log n}}_{\text{clipping}} .
\]
Indeed: the smoothing term is
$2S_n^{\psi}\log n\,C_{\mathrm{st}}\delta_n^{p}$ with
$S_n^{\psi}=n^{2(1-h)^2}$ and $\delta_n^p=n^{-hp}$, and
$2(1-h)^2-hp\le2(1-h)^2-h<0$ since $2(1-h)^2<\tfrac18<\tfrac34<h$; the
mismatch term is $S_n^{\psi}C_wC_{\log}n^{-(2-2b)}$ with exponent
$2(1-h)^2-(2-2b)$, negative exactly when $(1-h)^2<1-b$, i.e.\ (R4); the
clipping term is $2S_n^{\psi}\log n\,V_n(1+d'^{-2})
=O(n^{2(1-h)^2}\log n\cdot n^{-4(1-h)^2}/\log^2n)$, unconditionally
vanishing.
\end{proof}

\begin{remark}[Sharpness]
The mismatch rate is attained: the expansion of
$\log\|V_j\|^2$ has leading coefficient
$f'(s)s^{2h-1}\bigl[2hD'/D+2h\log s+1\bigr]$ evaluated at the window,
generically nonzero, so no per-row expansion order reaches the rate
$n^{-h}$ and the manuscript's original bias envelope at that rate fails for
non-constant profiles. The window refutations are machine-checked; the
lower-bound statement is derived by hand and registered.
\end{remark}

\section{Formalization cross-reference}

{\sloppy
Every \Wlean{} tag above is a theorem in the repository, axiom-checked within
\{\texttt{propext}, \texttt{Classical.\allowbreak{}choice}, \texttt{Quot.\allowbreak{}sound}\}, zero
\texttt{sorry}; the entry points are
\texttt{actualQ1\_\allowbreak\hspace{0pt}knownScaleH\_\allowbreak\hspace{0pt}fullChain\_\allowbreak\hspace{0pt}honestRate} (main theorem),
\texttt{honestRate\_\allowbreak{}drift\_\allowbreak{}tendsto\_\allowbreak{}zero} (Proposition~\ref{prop:drift-full}),
and the audit theorems listed in the companion note's provenance table.
Automated re-verification: \texttt{verification/\allowbreak{}final\_\allowbreak{}acceptance.\allowbreak{}py}.
\par}

\section*{References}\label{sec:refs}
\begin{enumerate}
\itemsep2pt
\item Source manuscript under audit: \texttt{19-AOS1825.pdf} (provided with
the project; not modified). Numbered results 3.1--3.5, 4.1--4.3, 8.1--8.5
and S.1.1--S.7.2 refer to it.
\item Companion notes: \texttt{paper\_revision/revised\_theorems.\{tex,pdf\}}
(the revised main theorem and the validation-status tables) and
\texttt{paper\_revision/theorem\_inventory.md} (per-item machine status with
Lean theorem names).
\item Full written chains (Chinese, authoritative for details not rendered
here): \texttt{direct\_proofs}/01--24. The per-item mapping is recorded in
\texttt{results/catalog.json}.
\item J.\ Bardet, M.\ Surgailis, \emph{Moment bounds and central limit
theorems for Gaussian subordinated arrays}, arXiv:1104.4732 --- Lemma 1 is
the only external lemma consumed by the Lean chain (as the hypothesis class
of \texttt{bardetSurgailis\_weighted\_evenMoment\_bound}); its scalar
rank-2 special case is used, as declared in \texttt{direct\_proofs}/19.
\item \emph{The functional Breuer--Major theorem}, arXiv:1808.02378 ---
cited for contrast only (\S1.1): the critical-regime proof here deliberately
avoids the Breuer--Major summability hypothesis, which fails at
criticality.
\end{enumerate}

\end{document}
